Galaxy NGC 2639 is morphologically identified as an Sa galaxy with measured
maximum rotational velocity $v_{\max}$ of 324\,km/s. After corrections for any
extinction, its apparent magnitude in $B$ is $m_B = 12.22$. It is customary to
measure a radius $R_{25}$ (in units of kpc) at which the galaxy's surface
brightness falls to $25\,\mathrm{mag}_B/\mathrm{arcsec}^2$. Spiral galaxies
tend to follow a typical relation:
\[ \log R_{25} = -0.249 M_B - 4.00, \]
where $M_B$ is the absolute magnitude in $B$. Apply the $B$-band Tully--Fisher
relation for Sa spirals
\[ M_B = -9.95 \log v_{\max} + 3.15 \quad (v_{\max}\ \text{in km/s}) \]
to calculate the mass of NGC 2639 out to $R_{25}$. If colour index of the sun
is $(m_{B_\odot} - m_{V_\odot}) = 0.64$, write the mass (of NGC 2639) in units
of solar mass $M_\odot$ and its luminosity $B$-band in unit of $L_\odot$.
